\chapter*{Приложение А}
\phantomsection\addcontentsline{toc}{chapter}{Приложение А. Примеры входных и выходных данных}
\begin{center}\textbf{Примеры входных и выходных данных}\end{center}

В таблице \ref{tab:app-examples} приведены примеры действий в интерфейсе. Входными данными служат название задачи и выбор записи или команды. Результатом становится изменение списка, статистики или сообщение пользователю.

\begingroup
\renewcommand{\thetable}{А.\arabic{table}}
\setcounter{table}{0}
\begin{table}[htbp]
\caption{Примеры действий пользователя и результатов}
\label{tab:app-examples}
\small\centering\renewcommand{\arraystretch}{1.28}
\begin{tabularx}{\textwidth}{|L{39mm}|Y|Y|}
\hline
Действие & Входные данные & Ожидаемый результат \\
\hline
Добавление & «Сдать практическую работу» & Новая строка в общем списке \\
\hline
Повторное чтение & Открытие списка на втором ПК & Та же строка появляется у второго клиента \\
\hline
Отметка выполнения & Выбрана добавленная задача & Перед названием показано [x] \\
\hline
Изменение названия & Выбрана задача «Подготовить курсовую» & Новое название сохраняется в базе \\
\hline
Поиск & Введено слово «курсовую» & Отображаются записи с совпадением \\
\hline
Фильтр & Выбрано «Выполнено» & Показаны только выполненные задачи \\
\hline
Статистика & Три задачи, одна выполнена & Доля выполнения составляет 33\% \\
\hline
Экспорт в TXT & Две задачи с разными отметками & В выбранном файле записаны обе задачи \\
\hline
Возврат в работу & Выбрана выполненная задача & Перед названием показано [ ] \\
\hline
Неверный ввод & Пустая строка или пробелы & Показана подсказка; запись не создаётся \\
\hline
Удаление & Выбрана существующая задача & Строка исчезает из общего списка \\
\hline
Отсутствие связи & Компьютер-сервер выключен & В строке состояния указано, что соединение отсутствует \\
\hline
\end{tabularx}
\end{table}
\endgroup



\chapter*{Приложение Б}
\phantomsection\addcontentsline{toc}{chapter}{Приложение Б. Порядок проверки на двух компьютерах}
\begin{center}\textbf{Порядок проверки на двух компьютерах}\end{center}

Для демонстрации на одном компьютере запускают \path{START.bat}: открывается сервер и два независимых окна клиента. Для проверки двух физических ПК на первом запускают \path{SERVER_ONLY.bat}, а на втором \path{CLIENT_OTHER_PC.bat} и вводят локальный IPv4-адрес первого компьютера.

При проверке на двух физических ПК оба компьютера должны находиться в одной локальной сети. Для определения IPv4-адреса на первом ПК выполняют команду \texttt{ipconfig}. Если второй компьютер не подключается, проверяют адрес сервера и правила брандмауэра для порта 8000.

После подключения на первом ПК создают задачу «Подготовить курсовую». На втором ПК проверяют появление записи после обновления списка. Затем исправляют название, меняют отметку выполнения и проверяют изменение на первом ПК. На каждом клиенте проверяют поиск, фильтр и показатели страницы «Статистика». Также нажимают «Сохранить TXT», выбирают файл и открывают его в Блокноте Windows. В завершение удаляют задачу. Для проверки сохранения создают новую запись, перезапускают сервер и убеждаются, что запись осталась в списке. Для проверки ошибки соединения временно останавливают сервер, ждут сообщения об отсутствии связи, затем запускают сервер снова и проверяют восстановление списка.

\vfill
\noindent Дата окончания курсового проекта: \rule{53mm}{0.4pt}\\[6mm]
\noindent Подпись автора: \rule{77mm}{0.4pt}
